¥chapter{プロセスとジョブ}

¥section*{プログラムを同時に実行}

最近のコンピュータは、パソコンのような小型のものでも、
ブラウザ、ゲーム、エディタなどを同時に起動して使うことができる。
これは同時にいくつものプログラムを実行できるということである。

このような機能をマルチタスク(multitask: 複数の仕事)という。
小型のコンピュータで、こんなことができるようになったのは、
本当に最近のことである。
実際、初期のコンピュータは、大型のものでも、
一度に一つのプログラムしか実行できなかった。

1970 年頃、 UNIX が開発された頃に、タイム・シェアリング
¥footnote{Time Sharing: 「時分割」とか訳す。
とても短い時間ごとにプログラムを切替えて、次々に少しずつ実行するという方法。}
などによって、
マルチタスクを実現することができるようになり、
ウィンドウをいくつも開いたり、
何人もの人で同じコンピュータを同時に使えるようになった。

走っているプログラムのことをプロセスと呼ぶが、
マルチタスクのシステムでは同時にプロセスがいくつも走っている。
今回は、このプロセスの見方、コントロールの仕方を紹介しよう。
プロセスがコントロールできるようになると、システムをスマートに使える
ようになるし、暴走したプログラムを止めたりすることもできるようになる。

¥section{プロセス}

¥subsection{プロセスの表示: ¥texttt{ps}}

 UNIX システムでは、 Kterm をいくつも開いたり、
時計(たとえば ¥texttt{xclock})を表示しながら、
 Mule を使ったりといった具合に、同時にいくつものプログラムを実行できる。

これは、システムが Kterm、¥texttt{xclock}、Mule などの
プログラムを同時に実行しているということだ。
実行中のコマンドのことを¥textbf{「プロセス」}という。

では、今、どんなプロセスが走っているのかを見てみよう。
走っているプロセスは ¥texttt{ps}
(¥textbf{P}rocess ¥textbf{S}tatus: プロセスの状態) コマンドで見ることができる。

¥begin{progls}
% ¥slbfl{ps}
  PID  TT  STAT      TIME COMMAND
  273  p0  Ss     0:00.07 -csh (tcsh)
  288  p0  R+     0:00.00 ps
%
¥end{progls}

なお、これは¥underline{実は全部ではない}
¥footnote{
表示されているのは、現在の端末から実行されているプロセスだけ。
少し難しい話なので、ここの部分はわからなくてもよい。}。

¥subsubsection{自分のプロセスを全部表示}

ある人のプロセスを全部表示するには「¥texttt{ps -U 「その人の ID」}」
と入力すればいい。

では、実際にやってみよう。
たとえば、 ¥texttt{ei00000} さんが、自分のプロセスを全部表示するには、
「¥texttt{ps -U ei00000}」である。

%¥begin{progls}
% ¥slbfl{ps -U ei00000}
%  PID  TT  STAT      TIME COMMAND
%  267  ??  Is     0:00.02 /bin/sh /usr/X11R6/lib/X11/xdm/Xsession
%  272  ??  S      0:00.08 twm
%  273  p0  Ss     0:00.07 -csh (tcsh)
%  289  p0  R+     0:00.00 ps -U ei00000
%
%¥end{progls}

¥begin{figure}[H]
¥centering{¥includegraphics[scale=0.8]{ps-u.eps}}
¥end{figure}

これが、走らせている全部のプロセスだ。

¥bigskip

この表示の見方を説明しよう。

まず、プロセスには、システムがコントロールするために付けた番号がある。
「¥texttt{PID}(プロセス ID)」というのがそれである。
システムの中には、同時に同じ番号のプロセスは存在しないようになっている。

次の列の「¥texttt{TT}」は、端末番号という。
実は、 Kterm のように、コマンドを打つことができるウィンドウを開くと、
システムがウィンドウごとに番号を割り振る。
それが端末番号で、これを見ると、どのウィンドウから実行したコマンドかがわかる
¥footnote{端末番号が「¥texttt{??}」になっているのは、 Kterm などからは
起動していないコマンドである。これらは、単に「¥texttt{ps}」と打つと表示されない。}。

次の「¥texttt{TIME}」は、これまでにプロセスが走ったトータルの時間。

最後の「¥texttt{COMMAND}」は、コマンドの名前である。

¥vspace*{6mm}

¥begin{BOXNOTE}
{¥large¥bf ☆練習☆}
¥begin{enumerate}
¥item まず、 X Window System で、いくつかのアプリケーションを起動してから、
¥texttt{ps}を実行し、どんなプロセスが走っているのかを見てみよう。
¥end{enumerate}
¥end{BOXNOTE}

¥subsubsection{すべてのプロセス}

 login した本人が動かしている以外にも、
システムそのものが動かしているプロセスがたくさんある。

マシンの中で動いているプロセスを全部表示するには、「¥texttt{ps -aux}」と
入力する。
これはたくさん出るので、次のようにページャー ¥texttt{jless} にパイプして
やるとよい。

¥begin{progls}
% ¥slbfl{ps -aux ¥PIPE jless}
¥end{progls}

すると、次のように表示される。
これが、現在実行中の全プロセスだ。
なお、 ¥texttt{USER} が「¥texttt{root}」や「¥texttt{daemon}」というのは、
システムが走らせているプロセスである。

¥begin{progls}
USER       PID %CPU %MEM   VSZ  RSS  TT  STAT STARTED       TIME COMMAND
ei00000    297  0.0  1.5   680  948  p0  RV    1Jan70    0:00.00 -csh (tcsh)
root         1  0.0  0.4   408  244  ??  Is    3:23AM    0:00.01 /sbin/init --
root         2  0.0  0.0     0   12  ??  DL    3:23AM    0:00.00  (pagedaemon)
root         3  0.0  0.0     0   12  ??  DL    3:23AM    0:00.00  (vmdaemon)
root         4  0.0  0.0     0   12  ??  DL    3:23AM    0:00.02  (update)
root        27  0.0  0.1   204   84  ??  Is    3:23AM    0:00.00 adjkerntz -i
root        90  0.0  0.8   208  488  ??  Ss    6:23PM    0:00.05 syslogd
daemon     100  0.0  0.9   176  576  ??  Is    6:23PM    0:00.00 portmap
root       104  0.0  0.9   180  544  ??  Ss    6:23PM    0:00.00 ypbind
root       114  0.0  0.2   212   88  ??  I     6:23PM    0:00.00 nfsiod -n 4
root       115  0.0  0.2   212   88  ??  I     6:23PM    0:00.00 nfsiod -n 4
root       116  0.0  0.2   212   88  ??  I     6:23PM    0:00.00 nfsiod -n 4
root       117  0.0  0.2   212   88  ??  I     6:23PM    0:00.00 nfsiod -n 4
root       132  0.0  1.0   240  604  ??  Is    6:23PM    0:00.02 inetd
root       135  0.0  0.8   364  516  ??  Is    6:23PM    0:00.01 cron
root       138  0.0  0.9   208  544  ??  Is    6:23PM    0:00.00 lpd
root       141  0.0  1.3   616  820  ??  Ss    6:23PM    0:00.00 sendmail: acce
root       164  0.0  0.7   168  420  ??  Ss    6:23PM    0:00.13 moused -p /dev
bin        224  0.0  0.9   540  584  ??  I     6:23PM    0:00.01 /usr/local/sbi
root       236  0.0  1.7   268 1064  ??  Is    6:23PM    0:00.01 /usr/X11R6/bin
root       244  0.8 13.7  3836 8640  ??  Ss    6:23PM    0:02.24 /usr/X11R6/bin
root       249  0.0  0.9   180  548  v0  Is+   6:23PM    0:00.01 /usr/libexec/g
:¥blsq
¥end{progls}

%¥newpage

このように ¥texttt{ps} は、次のように使う。
¥begin{coml}{6cm}
¥texttt{ps}

¥smallskip
¥texttt{ps -U 「ユーザの ID」}

¥smallskip
¥texttt{ps -aux}
¥end{coml}

¥newpage

¥subsection{CPU をよく使っているプロセス: ¥texttt{top}}

¥texttt{ps} 以外にも、プロセスを見る方法はある。
¥texttt{top} を使うと、 CPU を使っている順に、上位から順番にプロセスを
表示してくれる。
単に「¥texttt{top}」と打って見よう。
すると、下のようになる。
なお、 ¥texttt{top} を終了してこの状態から抜けるには「¥texttt{q}」と打つ。

¥begin{progls}
last pid:   294;  load averages:  0.04,  0.02,  0.00                   18:25:17
30 processes:  1 running, 29 sleeping
CPU states:  1.2% user,  0.0% nice,  0.0% system,  1.9% interrupt, 96.9% idle
Mem: 11M Active, 4956K Inact, 7324K Wired, 3790K Buf, 38M Free
Swap: 64M Total, 64K Used, 64M Free

  PID USERNAME PRI NICE SIZE    RES STATE    TIME   WCPU    CPU COMMAND
  244 root      2   0  3836K  8640K select   0:02  1.31%  1.30% XF86_SVGA
  294 ei00000  28   0   652K   860K RUN      0:00  0.40%  0.04% top
  279 ei00000  18   0   680K   984K pause    0:00  0.00%  0.00% tcsh
  273 ei00000  18   0   680K   940K pause    0:00  0.00%  0.00% tcsh
  135 root     18   0   364K   516K pause    0:00  0.00%  0.00% cron
   27 root     18   0   204K    84K pause    0:00  0.00%  0.00% adjkerntz
  252 root     10   0   352K  1600K wait     0:00  0.00%  0.00% xdm
  267 ei00000  10   0   492K   316K wait     0:00  0.00%  0.00% sh
    1 root     10   0   408K   244K wait     0:00  0.00%  0.00% init
  117 root     10   0   212K    88K nfsidl   0:00  0.00%  0.00% nfsiod
  114 root     10   0   212K    88K nfsidl   0:00  0.00%  0.00% nfsiod
  115 root     10   0   212K    88K nfsidl   0:00  0.00%  0.00% nfsiod
  116 root     10   0   212K    88K nfsidl   0:00  0.00%  0.00% nfsiod
  293 ei00000   3   0   600K   892K ttyin    0:00  0.00%  0.00% vi
  251 root      3   0   180K   548K ttyin    0:00  0.00%  0.00% getty
  250 root      3   0   180K   548K ttyin    0:00  0.00%  0.00% getty
  249 root      3   0   180K   548K ttyin    0:00  0.00%  0.00% getty
¥end{progls}

%¥newpage

¥subsection{プロセスを止める: ¥texttt{kill}}

コマンドの実行を止める方法を、ここまでにいくつか紹介してきた。
例えば ¥KEYTOP{Ctrl}+¥texttt{c} と入力する方法、また X Window System なら
Twm のメニューから「¥texttt{Delete}」を選択する方法などである。

ここではより細かい操作ができる方法を紹介しよう。

コマンドの実行を止めるのには、 ¥texttt{kill} を使ってプロセスを殺せばいい。
使い方は以下の通りである。

¥begin{coml}{4.5cm}
¥texttt{kill} 「PID」
¥end{coml}

では、実際にやってみよう。
まず、あらかじめ「¥texttt{xclock ¥&}」と打って、 ¥texttt{xclock} を
起動しておく。
¥begin{progls}
% ¥slbfl{xclock &}
[1] 301
%
¥end{progls}

¥newpage

次に、 ¥texttt{kill} を使うには、 PID を知る必要があるので、
次のように ¥texttt{ps} で ¥texttt{xclock} の PID を調べる。

¥begin{progls}
% ¥slbfl{ps -U ei00000}
  PID  TT  STAT      TIME COMMAND
  267  ??  Is     0:00.02 /bin/sh /usr/X11R6/lib/X11/xdm/Xsession
  272  ??  S      0:00.08 twm
  273  p0  Ss     0:00.07 -csh (tcsh)
  301  p0  S      0:00.03 xclock
  303  p0  R+     0:00.00 ps -U ei00000
%
¥end{progls}

すると、 ¥texttt{xclock} の PID は 301 だとわかる。
そこで「¥texttt{kill 301}」する。

¥begin{progls}
% ¥slbfl{kill 301}
%
¥end{progls}

すると、 ¥texttt{xclock} が消えたはずだ。
なお ¥texttt{kill}を実行すると、
「¥texttt{[1]    Done                          xclock}」
のような
メッセージが出ることがある。

では、念ため ¥texttt{ps} で確認してみよう。

¥begin{progls}
% ¥slbfl{ps -U ei00000}
  PID  TT  STAT      TIME COMMAND
  267  ??  Is     0:00.02 /bin/sh /usr/X11R6/lib/X11/xdm/Xsession
  272  ??  S      0:00.08 twm
  273  p0  Ss     0:00.07 -csh (tcsh)
  306  p0  R+     0:00.00 ps -U ei00000
¥end{progls}

確かに PID が 301 の ¥texttt{xclock} のプロセスが消えている。

ちなみに、このような方法でもプロセスが殺せない場合は、
強制的に殺す方法もある。
強制的に殺すには、「¥texttt{kill -KILL 「PID」}」と、
「¥texttt{-KILL}」オプションを
付けてやるといい。

¥vspace*{6mm}

¥begin{BOXNOTE}
{¥large¥bf ☆演習☆}
¥begin{enumerate}
¥item 実際に ¥texttt{kterm} や ¥texttt{xclock} などを起動しておいて、
 ¥texttt{kill} してみよう。
¥end{enumerate}
¥end{BOXNOTE}

¥newpage

¥section{ジョブ}

¥subsection{昔話: フォアグランド・ジョブとバックグランド・ジョブとは何か}

まず、イメージをつかみやすくするために、
昔のコンピュータの絵を見て欲しい。
次の図のように一台の大型のコンピュータがあり、
これに何台もの端末と呼ばれるものがつながっていた
¥footnote{図では主に「画面」と「キーボード」のついた端末を描いている。
しかし、実は最初の頃は左側に描いてるようにテレタイプライタといって、
「画面」ですらなくて、
代わりに「タイプライタの印字部」がついているようなものが使われていた。
そしてもっと前は・・・。}。

一方、今では X Window System で、一画面の中に何枚でも
気軽に Kterm のウィンドウを開くことができる。
 Kterm のウィンドウ一個一個は、昔の話では、
なんとそれぞれの端末に相当する
¥footnote{そういうわけで Kterm のようなソフトを
ターミナル・エミュレータ、あるいは仮想端末などと呼ぶことがある。}
のである。

¥begin{figure}[H]
¥centering{¥includegraphics[scale=0.4]{terminal.eps}}
¥caption{昔のコンピュータのイメージ}
¥end{figure}

%¥newpage

しかも、このそれぞれの端末は、文字しか表示できなくて、
画面の中に新しくウィンドウを開いたりすることもできなかった。
そして一度コマンドを打ったら、そのコマンドが実行し終わるまで
次のコマンドを入力することはできなかったのだ。

たとえば、大型コンピュータなどでも、計算¥footnote{有名な話では
流体計算、つまり大気、水、油などの流れる様子を計算する問題などがある。
天気予報や流体の研究などをするには、
よくスーパーコンピュータが使われていて、それでもすごい時間がかかった。}
をさせると何日もかかることが昔は(ものによっては現在でも)よくあった。
ある端末で時間のかかる計算をやらせるプログラムを書いて
実行させる。
すると、その端末は計算が終わるまで使えない・・・。
それではいくらなんでも不便だということで、
フォアグランド・ジョブ(foreground job: 画面の上でやる仕事)と
バックグランド・ジョブ(background job: 画面のバックでやる仕事)
というものが考えられた。

つまり、それぞれのユーザにとって、画面は一枚しかないので、
画面を使った仕事(フォアグランド・ジョブ)は一つしかできない。
しかしである、たとえば、計算だけしかプログラムなんて、
計算が終わって結果が出るまでは、
画面に何も出さなくていいはずじゃないかと。
そういう画面を使う必要のない計算は画面のバックでやってもらって、
画面は別な仕事¥footnote{メールを読み書きしたり、
新しい計算のプログラムを書いたり、
ゲームをしたり・・・。}に使えばいいじゃないかと。
そういう考え方である。

こうして、画面の表でやる仕事(フォアグランド・ジョブ)の他に、
画面の背後でも仕事(バックグランド・ジョブ)ができるように
OS が設計されるようになった。

¥begin{figure}[H]
¥centering{¥includegraphics[scale=1.0]{foreandback.eps}}
¥end{figure}

¥subsection{UNIX システムでのジョブのコントロール}

UNIX システムでは、ジョブというのは、だいたいプロセスのようなものと
思ってもらえばいい
¥footnote{厳密にはジョブとプロセスは異なっている。
例えば、単に「¥texttt{ls}」と打った場合、
プロセスもジョブも ¥texttt{ls} の 1 つだ。
しかし、「¥texttt{ls | jless}」と打った場合、
プロセスは ¥texttt{ls} と ¥texttt{jless} の 2 つ、
一方ジョブは ¥texttt{ls | jless} が 1 セットで 1 つになる。
ジョブというのは、このようにユーザが打ったコマンドの 1 まとまりのことである。}。
先ほどの説明通り、ジョブには、
フォアグランド(foreground: 表で実行)と、
バックグランド(background: 裏で実行)がある。
そしてフォアグランドで走らせられるジョブは一つだけ、
バックグラウンドでは複数のジョブを走らせることができる。

¥subsubsection{フォアグランド・ジョブの例}

UNIX システムでは、普通にコマンドを打つと、
それはフォアグランドのジョブになる。
これは、表で動いているジョブで、
 Kterm のような仮想端末において、
それぞれ同時に一つずつしか走らせられない。
つまり、普通に「¥texttt{ls}」とか打った場合、
 ¥texttt{ls} が終わるまで、次のコマンドは実行できないのだ。
とは言え、ほとんどの場合、一瞬で実行が終わるので気がつかないことだろう。

¥bigskip

では、敢えて時間がかかるようにして実行してみよう。
次のように「¥texttt{(sleep 30; ls)}」と入力してみよう。
「¥texttt{sleep 30}」というのは「30 秒休止」させるコマンドで、
このようにしてやると「30 秒休止してから ¥texttt{ls}を実行」する。

¥begin{progls}
% ¥slbfl{(sleep 30; ls)}
(30 秒間ストップ)
%
¥end{progls}

30 秒間は、この画面にどんなコマンドを打っても実行できないはずである。
この間に打ったコマンドは、30 秒たって ¥texttt{ls} が実行されてから、
順番に実行される。

¥subsubsection{バックグランド・ジョブの例}

バックグランド・ジョブは、一つの仮想端末でいくつも同時に走らせることができる。
UNIX システムでコマンドをバックグランドで走らせるには、
単にコマンドの最後に「¥texttt{¥&}」を付けるだけでいい。

¥bigskip

では、先ほどと同様のコマンドを今度はバックグランドで実行してみよう。

¥begin{progls}
% ¥slbfl{(sleep 30; ls) ¥&}
[1] 62749
%
¥end{progls}

30 秒後に ¥texttt{ls} が実行されるのは変わらない。
しかし、今回はその間にも他のコマンドが実行できるので、
試してみるといいだろう。
また、バックグランド・ジョブは同時にいくつも実行できるので、
時間差でいろいろなコマンドを実行してみるとおもしろいだろう。

¥begin{figure}[H]
¥centering{¥includegraphics[scale=0.8]{foreandback2.eps}}
¥end{figure}

¥subsubsection{フォアグランドをバックグランドに切り替え}

UNIX システムでは、
フォアグランド・ジョブとバックグランド・ジョブを切り替えることができる。
また、フォアグランド・ジョブの実行を「一時的に停止(サスペンド)」
することもできる。

具体的にやってみよう。
まず非常に時間のかかる(180 秒 = 3 分)フォアグランド・ジョブを実行してみよう。

¥begin{progls}
% ¥slbfl{(sleep 180; ls)}
¥end{progls}

フォアグランド・ジョブをバックグランド・ジョブに切り替えるには、
まず一時的にこれを停止(サスペンド)させる。
それには「¥KEYTOP{Ctrl}+z」と入力する。
すると、次のように一時的にジョブが停止し、画面に「¥texttt{{¥HAT}Z}」と表示される。

¥begin{progls}
% ¥slbfl{(sleep 180; ls)}
¥slbfl{^Z}
中断
%
¥end{progls}

次に「¥texttt{bg}」と打つと、一時停止中のジョブがバックグランドになる。
¥begin{progls}
% ¥slbfl{(sleep 180; ls)}
¥slbfl{^Z}
中断
% ¥slbfl{bg}
[1]    ( sleep 180; ls ) ¥&
%
¥end{progls}

このようにフォアグランド・ジョブは、
「¥KEYTOP{Ctrl}+¥texttt{z}」で一時停止(サスペンド)し、
「¥texttt{bg}」でバックグランド・ジョブに切り替えることができる。

¥begin{figure}[H]
¥centering{¥includegraphics[scale=0.8]{foretoback.eps}}
¥end{figure}

¥subsubsection{バックグランドをフォアグランドに切り替え}

次にバッググランド・ジョブを、
フォアグランド・ジョブに切り替える方法を紹介しよう。

さきほどと同様に非常に時間のかかる(180 秒 = 3 分)別なジョブを、
今度はバックグランドで実行してみよう。

¥begin{progls}
% ¥slbfl{(sleep 180; cal) ¥&}
[2] 62837
%
¥end{progls}

ここで、ジョブの一覧を見てみよう。
それには、以下のように ¥texttt{jobs} コマンドを使う。

¥begin{progls}
% ¥slbfl{jobs}
[1]    実行中です                    ( sleep 180; ls -CF )
[2]  + 実行中です                    ( sleep 180; cal )
¥end{progls}

¥texttt{jobs} コマンドを使うと、
現在実行中のバックグランド・ジョブの一覧を見ることができる。
先頭に書かれた「¥texttt{[1]}」や「¥texttt{[2]}」は
「¥textbf{ジョブ番号}」と呼ばれる。

これらのバックグランド・ジョブをフォアグランド・ジョブに
切り替えるには「¥texttt{fg  ¥%「ジョブ番号」}」と入力する。

では、次のようにしてジョブ番号が「1」の方をフォアグランドに切り替えてみよう。

¥begin{progls}
% ¥slbfl{fg %1}
( sleep 180; ls -CF )
¥end{progls}

このジョブがフォアグランドになって、実行が終わるまで他のコマンドは
実行できなくなったはずだ。

¥begin{figure}[H]
¥centering{¥includegraphics[scale=0.8]{backtofore.eps}}
¥end{figure}

¥subsubsection{X Window System アプリケーションの場合}

ここで思い出して欲しいのは、
 X Window System で新しくウィンドウを開くときのことである。
新しくウィンドウを開くコマンドを入力するときには、
最後に「¥texttt{¥&}」を付けていたが、
これは実はバックグランドでコマンドを実行していたのだ。

「¥texttt{¥&}」を付けないと、
フォアグランドでコマンドが実行されるから、
新しく開いたウィンドウを終了するまでは、
元のウィンドウでは何もできないのである。

これが昔のコンピュータの話と違うのは、
新しくウィンドウを開くコマンド
¥footnote{X Window System アプリケーションと呼ばれる。}
を実行することは、
新しい端末を用意することに相当するということである。

%
%¥bigskip
%
%少し説明がわかりにくかったかもしれないので、実際にやってみよう。
%
%まず、 Mule をフォアグランド(表)で走らせてみよう。
%単に ¥texttt{mule} と打とう。
%
%¥begin{progls}
%% ¥slbfl{mule}
%¥end{progls}
%
%こうすると、 Mule が起動する。
%
% Mule の方は、キーボードから値を入れたり、メニューを開いたりできる。
%しかし、 Mule を起動した ¥texttt{kterm} の方は、 ¥texttt{ls} とか
%打っても、実行されないはずだ。
%
%これは、今、 Mule がフォアグランドのジョブになっているということだ。
%フォアグランドのジョブは同時に 1 個しか走らないので、 ¥texttt{ls} とか
%打っても実行されないのだ。
%
%なお、 Mule の方を終了すれば、その間に ¥texttt{kterm} の方で打った
%コマンドが次々に実行される。
%
%¥bigskip
%
%では、次にバックグランドで Mule を起動してみよう。
%バックグランドで動かすには、
% ¥texttt{mule }¥verb+&+ と、最後に ¥verb+&+ を付ける。
%
%¥begin{progls}
%% ¥slbfl{mule &}
%[1] 544
%%
%¥end{progls}
%
%今度は、 Mule を起動した ¥texttt{kterm} の画面にも、コマンドが打てるはずだ。
%
%こうして、 X Window System でウィンドウを開く時には、
%バックグランドにしておけば、いくつも同時にコマンドを実行できる。
%
%ちなみに、 X Window System 以外のコマンドも ¥verb+&+ を最後に付けると、
%バックグランドで実行できる。
%これは、時間のかかるコマンドを実行させるときに便利だ。
%
%¥begin{figure}[H]
%%¥centering{¥includegraphics{bgfg.eps}}
%¥centering{¥includegraphics[scale=0.8]{foreback.eps}}
%¥end{figure}
%
%¥subsubsection{フォアグランドをバックグランドに変更}
%
%では、一度、フォアグランドで動かしてしまったジョブを、バックグランドに
%する方法を紹介しておこう。
%この話は少し難しいので、飛ばして読んでもかまわない。
%
%たとえば、 ¥texttt{kterm} から Mule をフォアグランドで実行してしまったが、
%途中でコマンドも実行したくなった場合、
% Mule をバックグランドに変更することができる。
%
%実際にやってみよう。
%まず、 Mule をフォアグランドで実行する。
%
%¥begin{progls}
%% ¥slbfl{mule}
%¥end{progls}
%
%ここで、 ¥texttt{kterm} の画面の方で、¥KEYTOP{Ctrl}+¥texttt{z} を入力すると、
%ジョブを一時停止(サスペンドと言う。終了ではないことに注意)できる。
%
%¥begin{progls}
%% ¥slbfl{mule}
%^Z
%Suspended
%%
%¥end{progls}
%
%今、 Mule は、一時停止状態で、何も入力できなくなっている。
%この状態で、 ¥texttt{bg} (バックグランドの略) と打つと、
%一時停止している Mule のジョブをバックグランドに変更して、
%再び入力できるようになる。
%¥begin{progls}
%% ¥slbfl{bg}
%[1]    mule &
%%
%¥end{progls}
%
%¥begin{figure}[H]
%%¥centering{¥includegraphics{suspend.eps}}
%¥centering{¥includegraphics[scale=0.8]{fore2back.eps}}
%¥end{figure}
%
¥section*{この章で紹介したコマンド}

¥begin{center}
¥begin{minipage}{13cm}
¥begin{itembox}[l]{プロセスとジョブ}
¥begin{description}
¥item[¥texttt{ps} :] プロセスを表示

使い方: ¥texttt{ps -U} 「ユーザ ID」、¥texttt{ps -aux}

¥item[¥texttt{top} :] CPU を使っている上位のプロセスを表示

使い方: ¥texttt{top}

¥item[¥texttt{kill} :] プロセスを殺す

使い方: ¥texttt{kill} 「PID」

¥item[¥texttt{jobs} :] ジョブの一覧を表示

¥item[¥texttt{¥KEYTOP{Ctrl}+¥texttt{z}} :] フォアグランド・ジョブを一時停止

¥item[¥texttt{bg} :] 一時停止中のジョブをバックグランドに

¥item[¥texttt{fg} :] バックグランド・ジョブをフォアグランド・ジョブに

使い方: ¥texttt{fg 「ジョブ番号」}
¥end{description}
¥end{itembox}
¥end{minipage}
¥end{center}

